\documentclass[12pt]{article}
\usepackage[utf8]{inputenc}
\usepackage[T1]{fontenc}
\usepackage{hyperref}
\title{Iyive Language Resources for K-12 Curriculum Development}
\author{Compiled by Claude Sonnet 5 (claude-sonnet-5) for Andr\'as Kornai}
\date{}
\begin{document}
\maketitle

\section{Introduction}
Iyive (also Yiive, Uive, Ndir; ISO 639-3 uiv) is a Central Tivoid (Southern Bantoid) language of Manyu Division, Southwest Region, Cameroon, and of Nigeria. Sources give about 2,000 speakers (1996) or fewer than 1,000, and describe it as severely endangered with little transmission to children. Very little material exists in the search results.

\section{Primary Educational Resources}
None found.

\section{Linguistic and Reference Materials}
SIL Cameroon archive, \emph{The Phonology Sketch of the Iyive Language}: \url{https://www.silcam.org/resources/archives/52746}. Overview: \url{https://en.wikipedia.org/wiki/Iyive_language}.

\section{Religious and Community Resources}
Joshua Project profile (Cameroon): \url{https://joshuaproject.net/index.php/people_groups/12302/CM}.

\section{Digital and Audio Resources}
The Endangered Languages Project holds sample entries for Iyive (listed under a Mesaka-related classification): \url{https://www.endangeredlanguages.com/lang/550/samples/9909?hl=es}.

\section{Curriculum Development Context}
Iyive is not known to be taught in schools. Sources disagree on whether it is written. A curriculum would have to begin with documentation and orthography work, building on the SIL phonology sketch.
\end{document}